\section{Haunted by the Emperor}

\begin{quotex}
Europe is haunted by the shadow of the Emperor. One senses his absence just as vividly as in former times one sensed his presence. Because the emptiness of the wound speaks, that which we miss knows how to make us sense it. \flright{\textsc{Valentin Tomberg}, \emph{Meditations on the Tarot}}

At the apex of everything, as myth and limit, the ideal of the \emph{Chakravartin}, ``the King of the World", the invisible emperor, whose strength is hidden, powerful, and unconditioned. \flright{\textsc{Julius Evola}, \emph{Pagan Imperialism}}

\end{quotex}
\paragraph{The Gelasian Dyarchy}
\textbf{Rene Guenon} wrote an essay describing how the traditional ordering of society was based on both Spiritual Authority and Temporal Power. \textbf{Ananda Coomaraswamy} followed that up with a similar essay describing its application in Indian political thinking. Probably less well known is that the notion of the two powers is firmly rooted also in the Western tradition. The following letter from Pope Gelasius I to Emperor Anastasius in 492 summarizes this teaching:

\begin{quotex}
There are two, august Emperor, by which this world is chiefly ruled, namely, the sacred authority (\emph{auctoritas sacrata}) of the priests and the royal power (\emph{regalis potestas}). Of these, that of the priests is weightier, since they have to render an account for even the kings of men in the divine judgment. You are also aware, most clement son, that while you are permitted honorably to rule over human kind, yet in divine matters you bend your neck devotedly to the bishops and await from them the means of your salvation. In the reception and proper disposition of the heavenly sacraments you recognize that you should be subordinate rather than superior to the religious order, and that in these things you depend on their judgment rather than wish to bend them to your will.

If the ministers of religion, recognizing the supremacy granted you from heaven in matters affecting the public order, obey your laws, lest otherwise they might obstruct the course of secular affairs by irrelevant considerations, with what readiness should you not yield them obedience to whom is assigned the dispensing of the sacred mysteries of religion? Accordingly, just as there is no slight danger m the case of the priests if they refrain from speaking when the service of the divinity requires, so there is no little risk for those who disdain—which God forbid—when they should obey. And if it is fitting that the hearts of the faithful should submit to all priests in general who properly administer divine affairs, how much the more is obedience due to the bishop of that see which the Most High ordained to be above all others, and which is consequently dutifully honored by the devotion of the whole Church. 

\end{quotex}
The obvious question, then, is why did not the Medieval civilization persist as long as other Traditional societies such as ancient Egypt, China, India, or even Rome. Clearly, the personal element entered into it, since the Pope and the Emperor often intruded into each other's domain. Then nationalism always threatened the spiritual unity and the political power of the Holy Roman Empire. That is why nationalism is anti-tradition.

Although Julius Evola was wildly incorrect in assuming the Emperor had dominion over the priests, his notion that the Emperor could also possess spiritual authority is worthy of consideration. Valentin Tomberg points out that the Pope and the Emperor are ``posts", not necessarily different individuals. \textbf{Vladimir Solovyov} tried to bring the idea of Tsar and High Priest back into Orthodoxy. However, while the role of the spiritual authority was explicitly defined, the limits of temporal power were not. For that, we can look at the reign of the Holy Roman Emperor Frederick II, even though the unity he created did not survive his death.

\paragraph{The Viking Emperor}
\textbf{Frederick II} was the descendant of the Normans who had invaded Sicily, although he also had German and Sicilian blood. His intelligence was huge, in stark contrast to the mediocre leaders of our own time. He was trained in scholastic and classical philosophy, but also studied Aristotle and the Islamic philosophers Avicenna and Averroes. Fluent in six languages, he was also a promoter of the arts, particularly poetry.

Frederick served as the model for \textbf{Dante}'s ideas on Monarchy. Thus, for him the State was as much part of the plan for salvation as was the Church. He also seems to have been the inspiration for Evola's description of the Emperor. Frederick claimed that he was the only true individual in the State, since he is the One, not a part of another (like his subjects).

As King of Sicily, he sought to unify the nation; this required several elements:

\begin{itemize}
\item community of blood 
\item community of speech 
\item community of festivals 
\item community of history 
\item unity of law 
\end{itemize}
This created a certain tension since it encouraged a nascent nationalism. The Church was less interested in blood ties, since it was a community of believers. Latin was preferred to the vernacular, and Church feasts to local festivals. This had to be carefully balanced, since it was the Church that guaranteed the Emperor's position as well as the Holy Roman Empire. This tension persists in our day.

\paragraph{The Rational Order}
Frederick's innovation was to base the State soundly on natural reason, so that it was not a matter of faith or belief. Hence it was purely natural, hence acceptable even to unbelievers. As such, the State was the resultant of three forces:

\begin{itemize}
\item \emph{Justitia} 
\item \emph{Necessitas} 
\item \emph{Providentia} 
\end{itemize}
To implement these ideas, Frederick created a separate hieratic order of intellectuals, modeled on the chivalric orders. This was not unlike the Prussian State under the Teutonic Order, since the constitution of the Prussian State was likewise based on a rational system. As a matter of fact, the Order of Teutonic Knights in Prussia and the Sicilian hierarchy mutually influenced each other.

\paragraph{Justice}
For Frederick, Justice is both the founder of, as well as founded in, the State. Kings and the State became necessary only because of the Fall. For Dante, the Roman Empire was symbolized by the Tree of Knowledge. The Emperor, then, would lead men back to the highest natural wisdom at the entrance to Paradise; then the Church would take up the task bring men back to that primeval state of innocence.

Adam was a law-breaker, hence the State and Justice became the remedy. Frederick implemented a new \emph{Book of Laws} for Sicily. Its purpose was Justice:

\begin{quotex}
daily to conceive new methods to reward the virtuous and to pulverize the vicious under repeated blows of punishment, when the old human laws under the changes wrought by time and circumstance no longer suffice to eradicate vice and to implant virtue. 

\end{quotex}
\paragraph{Necessity}
Basing himself on the teachings of Aristotle and Aristotle's Arabic commentators, Frederick saw the State and the King arising from natural necessity. Hence, there was no need for any other justification. Earlier leaders may have claimed that the State was instituted by God, or a god, but that could not compel belief. As belief in the gods waned, the state lost its support. However, Frederick claimed that the need for a ruler could be grasped by reason, or else the human race would have destroyed itself. Dante later elaborated on this idea in \emph{De Monarchia}, claiming:

\begin{quotex}
Thus Monarchy is necessary for the safety, for the advantage of the world. 

\end{quotex}
Of course, Necessitas encompassed more than that alone. It referred to the inevitability of things according to a natural law. Human life itself was subject to the law of Necessity. Rationality, then, results in the incorporation of the laws of nature into the positive law of the State.

\paragraph{Providence}
Frederick rejected the idea of a capricious God intervening miraculously into human affairs. Obviously, this is unlike the pagans, for whom the gods and goddesses were always interfering arbitrarily in world events. For Frederick, God's Providence was itself the Law, or the Logos as we would prefer to state it. That means that the natural order was itself the reflection of the Divine Order, and so Providence was Rational. The Scholastics had a similar definition:

\begin{quotex}
Providence is the Reason of a purposeful order of things. 

\end{quotex}
Moreover, it was the Emperor who was the mediator and interpreter of the Divine Plan.

\paragraph{Marriage Laws}
As an example of a specific application of Frederick's notion of Justice, we will consider the marriage laws that he implemented in Sicily. While not denying its sacramental value, Frederick regarded marriage as a necessity of nature for the preservation of the human race. Hence, the intention of those laws was improving the breed. Specifically, a Sicilian was forbidden to marry a non-Sicilian. This is his justification:

\begin{quotex}
It has often grieved us to see how the righteousness of our kingdom has suffered corruption from foreign manners by the mixture of different peoples. When the men of Sicily ally themselves with the daughters of foreigners, the purity of the race becomes besmirched, while evil and sensual weakness increase, the purity of the people is contaminated by the speech and by the habits of the others, and the seed of the stranger defiles the hearth of our faithful subjects. 

\end{quotex}
The Church was only concerned about the spiritual makeup of the marriage partners. That is, she follows a ``moral logic" (in Valentin Tomberg's terms), while the State follows an ``organic logic". This logic, moreover, is justified by the very nature of things.\footnote{We discussed the idea of the right of the state to determine its genetic makeup in \textit{Social Surgery} (\url{https://www.gornahoor.net/?p=8253}). Clearly, Frederick was centuries ahead of Bradley in this regard.}

\paragraph{Heresy Laws}
The Catholic Faith was declared the State Religion. Hence, heresy was not simply a crime against the Church, but was considered to be treason against the State. It was considered both blasphemy against God and an attack on the Emperor. Hence, heresy was a State matter and did not involve Church officials. In this, Frederick modeled these laws on Imperial Rome.

\paragraph{Multiculturalism}
The heresy laws did not apply to the Jews, Muslims, and Orthodox Greeks living in Sicily. Nevertheless, they were segregated from the Catholic Sicilians. Moreover, they were subject to their own laws. For example, the Jews were permitted to follow their own religious laws. They were even exempt from the laws against usury, which would have consequences years later. Although Jews and Muslims could not be oppressed, they were never full citizens. The same tolerance was not afforded to heretics or apostates.

The primary reference for this material is \emph{Frederick the Second} by Ernst Kantorowicz.



\flrightit{Posted on 2017-01-04 by Cologero }

\begin{center}* * *\end{center}

\begin{footnotesize}\begin{sffamily}



\texttt{Mark Citadel on 2017-01-18 at 12:05 said: }

A lot of food for thought here (Tomberg is one author I should certainly read more of). What would be your response to Maistre's critique of basing the state around `reason'? Is the kind of `reason' that Maistre criticized of a different order than Frederick's?

In your view, does the priestly caste have any judicial authority. If not for heresy, then for what, beyond intra-ecclesiastic discipline?


\hfill

\texttt{Reza on 2017-01-20 at 11:20 said: }

Just a minor correction: Avicenna was not an ``Arabic Philosopher". He was Iranian.

Otherwise a very interesting article.

[Thanks for bringing that to my attention … unfortunately, I was quoting the book]


\hfill

\texttt{Cologero on 2017-01-23 at 21:08 said: }

@Mark, I believe that Maistre had in mind attempts like Rousseau's to base the legitimacy of the state on pure reason; ultimately, there it can only be based on a spiritual authority. That does not imply that specific policies should be irrational; quite the contrary. Since you mentioned Tomberg, Tomberg identified orders of logic beyond formal logic: viz., organic logic and moral logic. That will be a topic here in the near future.


\hfill

\texttt{Mark Citadel on 2017-02-01 at 12:46 said: }

I look forward to it! What you have said makes sense, dividing the core of society (which is esoteric) from the core of policy, which is based on sound reason.

By the way, I recently moved my blog over to WordPress, so feel free to update the link on your blogroll whenever you get the time.

\url{http://citadelfoundations.wordpress.com}


\hfill


\end{sffamily}\end{footnotesize}
